%%%%%%%%%%%%%%%%%%%%%%%%%%%%%%%%%%%%%%%%%%%%%%%%%%%%%%%%%%%%%%%%%%%%%%
% Lớp 11 | Hàm số lượng giác và phương trình lượng giác | Các công thức lượng giác
% Dạng 1D1?3-1: Câu hỏi lý thuyết
% Số câu: 20  (sinh tự động bởi tach_cau.py; từ nay thầy sửa trực tiếp file này)
%%%%%%%%%%%%%%%%%%%%%%%%%%%%%%%%%%%%%%%%%%%%%%%%%%%%%%%%%%%%%%%%%%%%%%

%%%%%%%%%%%%%%%% TRẮC NGHIỆM 4 PHƯƠNG ÁN %%%%%%%%%%%%%%%%

% nguồn: Ok_Toan11_TracNghiem_Lan2
\begin{ex}%[1D1N3-1]%[ma:2c99cba7]
	Các đẳng thức nào sau đây đồng thời  xảy ra?
	\choice 
		{$\sin \alpha =\dfrac{\sqrt{2}}{3} $; $\cos\alpha  \dfrac{\sqrt{5}}{3} $}
		{$\sin \alpha =\dfrac{\sqrt{2}}{5} $; $\cos\alpha \dfrac{\sqrt{5}}{5}  $}
		{\True $\sin \alpha =\dfrac{4}{5} $; $\cos\alpha \dfrac{-3}{5}  $}
		{$\sin \alpha =\dfrac{\sqrt{3}}{4} $; $\cos\alpha \dfrac{1}{4}  $}
	\loigiai{ 
		Áp dụng công thức $\sin  ^{2} \alpha +\cos^2 \alpha =1$.\\
	Với  $\sin \alpha =\dfrac{4}{5} $; $\cos\alpha \dfrac{-3}{5} \Rightarrow\left(\dfrac{4}{5}\right)^2+ \left(\dfrac{-3}{5}\right)^2=1$.}
\end{ex}

% nguồn: Ok_Toan11_TracNghiem_Lan2
\begin{ex}%[1D1N3-1]%[ma:d792a179]
	Trong  các công thức sau, công thức nào đúng?
	\choice 
		{$\sin (a+b)=\sin  a \cdot \cos b-\cos a \cdot \sin b$}
		{$\cos (a+b)=\cos  a \cdot \cos b+\sin a \cdot \sin b$}
		{$\sin (a-b)=\sin a \cdot \cos b+\cos  a \cdot \sin b$}
		{\True $\cos (a-b)=\cos a \cdot \cos b+\sin  a \cdot \sin b$}
	\loigiai{ 
		Theo công thức cộng  ta có
		\begin{itemize}
			\item $\cos (a \pm b)=\cos a \cdot \cos b \mp \sin a  \cdot \sin b$
			\item $\sin  (a \pm b)=\sin a \cdot \cos b \pm \cos a \cdot \sin b$
	\end{itemize} }
\end{ex}

% nguồn: Ok_Toan11_TracNghiem_Lan2
\begin{ex}%[1D1N3-1]%[ma:95a68279]
	Cho $a$ là số  thực bất kỳ. khẳng định nào sau đây đúng?
	\choice 
		 {$\sin a+\cos a=1$}
		{$\sin ^{3} a+\cos ^{3} a=1$} 
		{$\sin ^{4}  a+\cos ^{4} a=1$}
		{\True $\sin ^{2} a+\cos ^{2}  a=1$}
	 \loigiai{
	Ta  có $\sin ^{2} a+\cos ^{2} a=1$; $\forall a\in \mathbb{R}$ nên $\sin ^{2} a+\cos ^{2} a=1$ đúng.}
\end{ex}

% nguồn: Ok_Toan11_TracNghiem_Lan2
\begin{ex}%[1D1N3-1]%[ma:0a2689c0]
	Trong các hệ thức sau đây, hệ thức nào  đúng?
	 \choice
		{\True $\sin x^2 +\cos x^2 =1$} 
		{$\sin x+\cos  x=1$}
		{$\sin 2x+\cos  2x=1$}
		{$\sin ^{2} x+\cos  x^2 =1$}
	 \loigiai{
	Theo công  thức $\sin x^2 +\cos x^2 =1$.}
\end{ex}

% nguồn: Ok_Toan11_TracNghiem_Lan2
\begin{ex}%[1D1N3-1]%[ma:473607ef]
	Mệnh  đề nào sau đây đúng?
	 \choice
		{$\tan (x-y)=\dfrac{\tan x+\tan y}{\tan  x \tan y}$}
		{\True $\tan (x-y)=\dfrac{\tan x-\tan y}{1+\tan x \tan  y}$}
		{$\tan (x-y)=\dfrac{\tan x-\tan y}{1-\tan x \tan  y}$}
		{$\tan  (x-y)=\dfrac{\tan x-\tan y}{\tan x \tan y}$}
	\loigiai{ 
	Ta có $\tan (x-y)=\dfrac{\tan x-\tan y}{1+\tan x  \tan y}$.}
\end{ex}

% nguồn: Ok_Toan11_TracNghiem_Lan2
\begin{ex}%[1D1N3-1]%[ma:63aac56a]
	Hãy chọn kết quả \textbf{sai} trong các kết quả sau đây 
	\choice 
		 {$\cos (-\alpha)=\cos \alpha $}
		{\True $\sin  (\pi +\alpha)=\sin \alpha $}
		{$\tan (\pi -\alpha)=-\tan \alpha  $}
		{$\cot \left(\dfrac{\pi}{2}  -\alpha \right)=\tan \alpha $}
	 \loigiai{
	Ta có $\sin (\pi +\alpha)=-\sin \alpha $. Do đo $\sin (\pi +\alpha)=\sin \alpha  $ là sai.}
\end{ex}

% nguồn: Ok_Toan11_TracNghiem_Lan2
\begin{ex}%[1D1N3-1]%[ma:9c0bfa72]
	Phát biểu  nào sau đây đúng?
	\choice 
		{$\tan (\alpha+\beta)=\dfrac{\tan \alpha-\tan \beta}{1+\tan  \alpha \cdot \tan \beta}$}
		{$\tan (\alpha+\beta)=\dfrac{1+\tan \alpha \cdot \tan  \beta}{\tan \alpha-\tan \beta}$}
		{\True $\tan  (\alpha-\beta)=\dfrac{\tan \alpha-\tan \beta}{1+\tan \alpha \cdot \tan \beta}$}
		{$\tan (\alpha-\beta)=\dfrac{1-\tan  \alpha \cdot \tan \beta}{\tan \alpha+\tan \beta}$}
	\loigiai{ 
	Theo công thức cộng ta có $\tan (\alpha \pm  \beta)=\dfrac{\tan \alpha \pm \tan \beta}{1 \mp \tan \alpha \cdot \tan \beta}$.}
\end{ex}

% nguồn: Ok_Toan11_TracNghiem_Lan2
\begin{ex}%[1D1N3-1]%[ma:8e12beaf]
	Chọn  khẳng định \textbf{sai} trong các khẳng định sau
	\choice 
		{\True $\cos  (a+b)=\cos a \cos b+\sin a \sin b$}
		{$\sin (a+b)=\sin a \cos  b+\cos a \sin b$}
		{$\sin (a-b)=\sin a \cos b-\cos a \sin  b$}
		{$\cos 2 a=1-2 \sin  ^{2} a$}
	 \loigiai{
	Ta có công thức đúng là $\cos (a+b)=\cos a \cos b-\sin  a \sin b$.}
\end{ex}

% nguồn: Ok_Toan11_TracNghiem_Lan2
\begin{ex}%[1D1N3-1]%[ma:444c8001]
	Trong các công thức  sau, công thức nào đúng?
	\choice 
		{\True $\sin (a-b)=\sin a  \cdot \cos b-\cos a \cdot \sin b$}
		{$\cos (a-b)=\cos a \cdot \cos b-\sin  a \cdot \sin b$}
		{$\sin (a+b)=\sin a  \cdot \cos b-\cos a \cdot \sin b$}
		{$\cos (a+b)=\cos a  \cdot \cos b+\sin a \cdot \sin b$}
	 \loigiai{
	 Công thức cộng là $\sin (a-b)=\sin a \cdot \cos b-\cos a \cdot \sin b$.}
\end{ex}

% nguồn: Ok_Toan11_TracNghiem_Lan2
\begin{ex}%[1D1N3-1]%[ma:999dc665]
	Biểu thức $\sin x \cos y-\cos x \sin y$ bằng 
	\choice 
		{$\cos  (x-y)$}
		{$\cos (x+y)$} 
		 {\True $\sin (x-y)$}
		{$\sin  (y-x)$}
	\loigiai{Ta có công thức $\sin  (x-y)=\sin x \cos y-\cos x \sin y$.}
\end{ex}

% nguồn: Ok_Toan11_TracNghiem_Lan2
\begin{ex}%[1D1N3-1]%[ma:98813c7f]
	Trong các khẳng định  sau, khẳng định nào \textbf{sai}?
	 \choice
		{$\sin a-\sin b=2  \cos \dfrac{a+b}{2} \sin \dfrac{a-b}{2}$}
		{\True $\cos (a-b)=\cos a \cos b-\sin a \sin b$} 
		{$\sin (a-b)=\sin  a \cos b-\cos a \sin b$}
		{$2 \cos a  \cos b=\cos (a-b)+\cos (a+b)$}
	 \loigiai{
	 Vì $\cos (a-b)=\cos a \cos b+\sin a \sin b$.}
\end{ex}

% nguồn: Ok_TracNghiemLop11_Lan1
\begin{ex}%[1D1N3-1]%[ma:9019ad80]
	Trong các công thức sau, công  thức nào đúng?
	 \choice
		{$\cos (a-b)=\cos a\sin b+\sin  a\cos b$}
		{\True $\sin (a-b) =  \sin a\cos b - \cos a\sin b$}
		{$\sin (a+b) = \sin a\cos b - \cos a\sin b$} 
		{$\cos (a+b)=\cos a\cos b+\sin a\sin  b$}
	 \loigiai{
		Công thức đúng: $\sin (a-b)=\sin a\cos  b-\cos a\sin b$.
	} 
\end{ex}

% nguồn: Ok_TracNghiemLop11_Lan1
\begin{ex}%[1D1N3-1]%[ma:7454baac]
	Công thức nào sau đây  đúng?
	 \choice
		{$\cos(2\beta)=2\sin^2\beta-1$} 
		 {\True $\cos(2\beta)=\cos^2\beta-\sin^2\beta$}
		{$\cos(2\beta)=\sin^2\beta-\cos^2\beta$} 
		{$\cos(2\beta)=1-2\cos^2\beta$} 
	 \loigiai{
		Dùng  công thức nhân đôi
		\begin{eqnarray*}
			 \cos(2\beta)&=&2\cos^2(\beta)-1\\
			&=&1-2\sin^2(\beta)\\ 
			 &=&\cos^2(\beta)-\sin^2(\beta).
		\end{eqnarray*}
	 }
\end{ex}

% nguồn: Ok_TracNghiemLop11_Lan1
\begin{ex}%[1D1N3-1]%[ma:6f4b7d0f]
	Trong các đẳng  thức sau, đẳng thức nào \textbf{sai}?
	\choice 
		{$\sin  a-\sin b=2 \cos \dfrac{a+b}{2}\sin \dfrac{a-b}{2}$}
		{$\cos a+\cos b=2 \cos \dfrac{a+b}{2}\cos  \dfrac{a-b}{2}$}
		{\True $\cos a-\cos b=2 \sin \dfrac{a+b}{2}\sin \dfrac{a-b}{2}$} 
		{$\sin  a+\sin b=2 \sin \dfrac{a+b}{2}\cos \dfrac{a-b}{2}$}
	\loigiai{ 
		Ta có $\cos a-\cos b=-2 \sin  \dfrac{a+b}{2}\sin \dfrac{a-b}{2}$.
	 }
\end{ex}

% nguồn: Ok_TracNghiemLop11_Lan1
\begin{ex}%[1D1N3-1]%[ma:f6040a06]
	Trong các mệnh đề sau, mệnh đề nào là  mệnh đề \textbf{sai}?
	 \choice
		{$\sin a+\sin b=2\sin \dfrac{a+b}{2}\cos \dfrac{a-b}{2}$} 
		{\True $\cos a-\cos b=2\sin \dfrac{a+b}{2}\sin  \dfrac{a-b}{2}$}
		{$\cos a+\cos b=2\cos \dfrac{a+b}{2}\cos  \dfrac{a-b}{2}$}
		{$\sin  a-\sin b=2\cos \dfrac{a+b}{2}\sin \dfrac{a-b}{2}$}
	 \loigiai{
	Công thức đúng là $\cos a-\cos b=-2\sin \dfrac{a+b}{2}\sin  \dfrac{a-b}{2}$.}
\end{ex}

% nguồn: Ok_Toan11_TracNghiem_Lan2
\begin{ex}%[1D1H3-1]%[ma:ecba060f]
	Cho hai góc $\alpha$ và $\beta$ thỏa mãn  $\sin \alpha=\dfrac{3}{5}$, $\left(\dfrac{\pi}{2}<\alpha<\pi\right)$ và $\cos \beta=\dfrac{12}{13}$, $\left(0<\beta<\dfrac{\pi}{2}\right)$. Giá trị của $\sin (\alpha-\beta)$ là
	 \choice
		{$-\dfrac{56}{65}$} 
		{\True  $\dfrac{56}{65}$}
		 {$\dfrac{16}{65}$}
		 {$-\dfrac{16}{65}$}
	\loigiai{ 
		Ta có $\dfrac{\pi}{2}<\alpha<\pi$ nên $\cos \alpha<0  \Rightarrow \cos \alpha=-\sqrt{1-\sin ^{2} \alpha}=-\dfrac{4}{5}$.\\
		Lại có $0<\beta<\dfrac{\pi}{2}$ nên $\sin \beta>0 \Rightarrow  \sin \beta=\sqrt{1-\cos ^{2} \alpha}=\dfrac{5}{13}$.\\
	Vậy $\sin (\alpha-\beta)=\sin \alpha \cos \beta-\cos \alpha \sin \beta=\dfrac{3}{5} \cdot \dfrac{12}{13}-\left(\dfrac{-4}{5}\right)  \cdot \dfrac{5}{13}=\dfrac{56}{65}$.}
\end{ex}

% nguồn: Ok_Toan11_TracNghiem_Lan2
\begin{ex}%[1D1H3-1]%[ma:068356da]
	Cho $\sin \alpha=\dfrac{2 \sqrt{5}}{5}$ với $0<\alpha<\dfrac{\pi}{2}$. Biết giá  trị của $\cos \left(\alpha+\dfrac{\pi}{3}\right)=\dfrac{a \sqrt{5}-b \sqrt{15}}{10}$ với $a$, $b\in \mathbb{N}$ và $(a, b)=1$. Tính $a+b$.
	 \choice
		{$4$ } 
		 {$10$ }
		{$7$ } 
		 {\True$3$ }
	 \loigiai{
		Do  $0< \alpha< \dfrac{\pi}{2}$ nên $\cos \alpha >0$.\\
		Ta có $$\sin ^{2} \alpha+\cos  ^{2} \alpha=1 \Leftrightarrow \cos ^{2} \alpha=\dfrac{1}{5} \Leftrightarrow \cos \alpha=\dfrac{\sqrt{5}}{5}.$$
		Ta có $$\cos \left(\alpha+\dfrac{\pi}{3}\right)=\dfrac{1}{2} \cos \alpha-\dfrac{\sqrt{3}}{2} \sin  \alpha=\dfrac{1}{2} \cdot \dfrac{\sqrt{5}}{5}-\dfrac{\sqrt{3}}{2} \cdot \dfrac{2 \sqrt{5}}{5}=\dfrac{\sqrt{5}}{10}-\dfrac{2 \sqrt{15}}{10}=\dfrac{\sqrt{5}-2 \sqrt{15}}{10}.$$
	Suy ra $a=1$, $b=2 \Rightarrow a+b=3$.} 
\end{ex}

% nguồn: Ok_Toan11_TracNghiem_Lan2
\begin{ex}%[1D1H3-1]%[ma:38b3a55f]
	Tính giá trị  $\cos \left(\alpha-\dfrac{\pi}{6}\right)$ biết $\sin \alpha=\dfrac{1}{3}$, $ \dfrac{\pi}{2}<\alpha<\pi$.
	\choice 
		{$-\dfrac{2  \sqrt{2}}{3}$}
		{$-\dfrac{1+2  \sqrt{6}}{6}$}
		{\True  $\dfrac{1-2 \sqrt{6}}{6}$}
		{$\dfrac{1+2 \sqrt{6}}{6}$} 
	 \loigiai{
		Vì $\sin  \alpha=\dfrac{1}{3}$, $\dfrac{\pi}{2}<\alpha<\pi$ nên $\cos \alpha=-\dfrac{2 \sqrt{2}}{3}$.\\
	Do đó $\cos \left(\alpha-\dfrac{\pi}{6}\right)=\cos \alpha \cdot \cos  \dfrac{\pi}{6}+\sin \alpha \cdot \sin \dfrac{\pi}{6}=-\dfrac{2 \sqrt{2}}{3} \cdot \dfrac{\sqrt{3}}{2}+\dfrac{1}{3} \cdot \dfrac{1}{2}=\dfrac{1-2 \sqrt{6}}{6}$.}
\end{ex}

% nguồn: Ok_Toan11_TracNghiem_Lan2
\begin{ex}%[1D1H3-1]%[ma:a3fcebd3]
	Rút gọn biểu thức: $\sin \left(a-17^{\circ}\right) \cdot \cos \left(a+13^{\circ}\right)-\sin  \left(a+13^{\circ}\right) \cdot \cos \left(a-17^{\circ}\right)$, ta được
	\choice 
		{$\sin 2 a$} 
		{$\cos 2  a$}
		{\True $-\dfrac{1}{2}$} 
		 {$\dfrac{1}{2}$}
	\loigiai{ 
		Ta có\begin{align*}
		\sin \left(a-17^{\circ}\right) \cdot \cos  \left(a+13^{\circ}\right)-\sin \left(a+13^{\circ}\right) \cdot \cos \left(a-17^{\circ}\right)&=\sin \left[\left(a-17^{\circ}\right)-\left(a+13^{\circ}\right)\right]\\
		&=\sin \left(-30^{\circ}\right)=-\dfrac{1}{2}. 
\end{align*} }
\end{ex}

% nguồn: Ok_Toan11_TracNghiem_Lan2
\begin{ex}%[1D1H3-1]%[ma:def84681]
	Cho góc lượng giác $\alpha$ thoả mãn $\left(\dfrac{\pi}{2}<\alpha<\pi\right)$. Xét  dấu $\sin \left(\alpha+\dfrac{\pi}{2}\right)$ và $\tan (-\alpha)$. Chọn kết quả đúng.
	\choice 
		 {$\heva{&\sin \left(\alpha+\dfrac{\pi}{2}\right)<0 \\ &\tan (-\alpha)<0}$}
		{$\heva{&\sin  \left(\alpha+\dfrac{\pi}{2}\right)>0 \\ &\tan (-\alpha)<0}$}
		{\True $\heva{&\sin  \left(\alpha+\dfrac{\pi}{2}\right)<0 \\& \tan (-\alpha)>0}$}
		{$\heva{&\sin \left(\alpha+\dfrac{\pi}{2}\right)>0 \\& \tan (-\alpha)>0}$} 
	\loigiai{ 
	Ta có $\dfrac{\pi}{2}<\alpha<\pi \Rightarrow\heva{&\pi<\alpha+\dfrac{\pi}{2}<\dfrac{3 \pi}{2} \\ &-\pi<-\alpha<-\dfrac{\pi}{2}} \Rightarrow\heva{&\sin \left(\alpha+\dfrac{\pi}{2}\right)<0  \\& \tan (-\alpha)>0.}$}
\end{ex}
